% Title: Sourdough Pizza Dough
% Description: A simple overnight pizza dough made from sourdough starter.
% Author: Wesley Jones
% Date: January 2025

\documentclass{article}

\usepackage[nonumber,noindex,nocontents]{cuisine}
\usepackage[margin=1in]{geometry}
\usepackage{endnotes}

% Set widths
\RecipeWidths  
  {1\textwidth}    %{Total recipe width}             
  {0.1\textwidth}  %{Number of servings width}                   
  {0.05\textwidth} %{Step number width}             
  {0.15\textwidth} %{Ingredient width} 
  {0.1\textwidth}  %{Quantity width} 
  {0.1\textwidth}  %{Units width}


% End notes instead of footnotes
\let\footnote=\endnote


\begin{document}

\begin{recipe}{Sourdough Pizza Dough}{\textit{Serves: 4}}{\textit{Preparation Time: 10 minutes + 1 day}}
  \freeform
    You will be able to make the dough with very little of your own interaction.
    The key  to this working is time!
    Pay attention to the timings indicated in the first two steps of this recipe --- they occur the \textit{night before} you intend to eat your pizza.
    The timings are \textit{suggestions}. An overnight is around 8 hours --- and the day time is over 8 hours.

  \freeform
    \\
    \textbf{\textit{The Night Before}}

  \Ingredient{40 grams all-purpose flour}
  \Ingredient{40 grams water}
  \Ingredient{40 grams starter}
  \textbf{Make the levain.}
  Comebine ingredients into mixing bowl the \textit{night before} you intend to make pizza.
  Stir to combine, cover and leave sit on counter overnight (or around 8 hours).
  
  \Ingredient{30 grams all-purpose flour}
  \Ingredient{30 grams water}
  \textbf{Refresh the starter.}
  Combine flour and water into starter.
  Stir to combine well and cover.
  Leave sit on counter overnight or for around 8 hours.

  \freeform
    \\
    \textbf{\textit{The Morning Of}}

  \newstep
    Place starter from Step 2 back into fridge. \textit{Until next time lil' guy!}

  \Ingredient{200 grams water}
  \Ingredient{10 grams olive oil}
  \Ingredient{330 grams flour, plus additional}
  \Ingredient{10 grams salt}
  \textbf{Make the dough.}
  Combine water and oil
    \footnote{This order ensures that the oil does not sticky up the flour and create a hard to combine dough}
  and then flour and salt into the levain from Step 1.
  Use mixer or mix by hand to integrate into a dough.
  The result will be very tacky and slightly slimey.
  You can add small quantities of flour (use a Tbsp to help regulate) if the dough is soupy.
  There is a good chance that adding flour will prove unnecessary, if you are uncertain, skip it.
  \\
  Allow dough to rest on the counter, covered, until pizza time.
 
  \freeform
    \\
    \textbf{\textit{During the Afternoon}}\footnote{
      Once the dough is made it can fridged at \textit{any} point for around 2 days.
      Don't panic if your plans change or if you want to work ahead, simply chuck the dough in the fride.
    }
  \newstep
    \textbf{Stretch \& fold the dough.}
    Several times over the course of the last 4 hours of dough resting, fold the dough over on itself.
    This can be done by simply using a spatula to bring part of the dough up and \textit{flipping} back onto the rest.
    Several of these flips should do the trick.
    2-3 interactions (once per hour) are enough.
    This step helps, but your pizza will turn out regardless.

  \freeform
    \\
    \textbf{\textit{Pizza Time}}\footnote{
      \textbf{Tools:} A pizza steel, a stone, a whole dedicated pizza oven --- or just a pan and conventional oven --- any will do. This dough is tough and forgiving. If you have variating tools you'll want to adjust to fit those tools' needs and functions.
    }
  \newstep
    \textbf{Make pizza.}
    Preheat to as hot as it will go (ideally at least 500\0).
    Flip the dough out onto a floured surface --- a good amount of flour (several Tbsps.).
    Flatten the dough out to an even round disc.
    Cover and let rest for 30 minutes.
    \\
    Coat a convential pizza pan with olive oil.
    Stretch the dough out carefully to fill a 16-inch round pan.
    You can use any level of skill, but the goal is an even layer accross the pan, no holes, with a small ridge for a larger crust (if desired).
    Top with toppings, bake on lowest rack for around 20 minutes.
    Check the bottom for done-ness.
    On a typical sheet or round pan, it will have a soft golden brown and lift up cleanly to support the whole pizza.
    Do not be afraid to take a pizza spatula and lift up to look and get the feel.
    Typically your toppings will begin to brown and the top crust will show good color at the same time as the bottom is done.\footnote{
      If you've added a lot of toppings, feel free to use the broil function to finish the toppings and cheese.
    }

\end{recipe}

\theendnotes

\end{document}
